\documentclass[11pt]{article}
\usepackage[margin=1in]{geometry}
\usepackage{amsmath,amssymb,amsthm}
\usepackage[utf8]{inputenc}
\usepackage{listings}
\usepackage{hyperref}
\usepackage{xcolor}
\usepackage{textcomp}

\definecolor{leanblue}{rgb}{0.2,0.4,0.7}

\lstdefinelanguage{Lean}{
  keywords={theorem,lemma,def,axiom,structure,class,instance,where,by,sorry,simp,exact,intro,apply,have,let,fun,match,if,then,else,import,open,namespace,end,section,variable,inductive,noncomputable,abbrev,deriving},
  keywordstyle=\color{leanblue}\bfseries,
  comment=[l]{--},
  morecomment=[s]{/-}{-/},
  string=[b]",
  basicstyle=\ttfamily\small,
  breaklines=true,
  extendedchars=true,
  inputencoding=utf8,
  literate=
    {≥}{{$\geq$}}1
    {≤}{{$\leq$}}1
    {→}{{$\to$}}1
    {←}{{$\leftarrow$}}1
    {∧}{{$\wedge$}}1
    {∨}{{$\vee$}}1
    {¬}{{$\neg$}}1
    {∀}{{$\forall$}}1
    {∃}{{$\exists$}}1
    {⊆}{{$\subseteq$}}1
    {⊂}{{$\subset$}}1
    {∈}{{$\in$}}1
    {∉}{{$\notin$}}1
    {×}{{$\times$}}1
    {⊗}{{$\otimes$}}1
    {▸}{{$\blacktriangleright$}}1
    {⟨}{{$\langle$}}1
    {⟩}{{$\rangle$}}1
    {λ}{{$\lambda$}}1
    {α}{{$\alpha$}}1
    {β}{{$\beta$}}1
    {σ}{{$\sigma$}}1
    {θ}{{$\theta$}}1
    {ε}{{$\varepsilon$}}1
    {μ}{{$\mu$}}1
    {π}{{$\pi$}}1
    {Σ}{{$\Sigma$}}1
    {Π}{{$\Pi$}}1
    {▶}{{$\triangleright$}}1
    {⊔}{{$\sqcup$}}1
    {⊓}{{$\sqcap$}}1
    {⊥}{{$\bot$}}1
    {⊤}{{$\top$}}1
    {∅}{{$\emptyset$}}1
    {∪}{{$\cup$}}1
    {∩}{{$\cap$}}1
    {≠}{{$\neq$}}1
    {ℕ}{{$\mathbb{N}$}}1
    {₀}{{$_0$}}1
    {₁}{{$_1$}}1
    {₂}{{$_2$}}1
    {|>}{{$|{\!>}$}}2
    {·}{{$\cdot$}}1
    {—}{{---}}1
    {∘}{{$\circ$}}1
    {√}{{$\sqrt{}$}}1
    {∝}{{$\propto$}}1
    {≈}{{$\approx$}}1
    {═}{{=}}1
    {⟹}{{$\Longrightarrow$}}1,
}

\newtheorem{theorem}{Theorem}
\newtheorem{lemma}[theorem]{Lemma}
\newtheorem{definition}[theorem]{Definition}
\newtheorem{axiom_stmt}[theorem]{Axiom}

\title{Formal Verification: CKKS Approximate Arithmetic FHE}
\author{Zach Kelling \and Woo Bin}
\date{December 2025}

\begin{document}
\maketitle

\begin{abstract}
We formalize CKKS, the approximate arithmetic FHE scheme used for ML inference on encrypted data. CKKS supports homomorphic addition and multiplication on approximate real numbers, with level management for multiplicative depth and SIMD batching for throughput.
\end{abstract}

\section{Introduction}
CKKS enables private ML inference by computing on encrypted floating-point vectors. With SIMD batching, a single ciphertext operation processes thousands of values in parallel. Combined with GPU acceleration, this enables practical encrypted transformer inference on Lux.

\section{Definitions}
\begin{definition}[Params]
CKKS parameters: polynomial degree $N$, scale bits, multiplicative levels.
\end{definition}

\begin{definition}[CKKSCiphertext]
An encrypted vector with remaining level, scale factor, and validity.
\end{definition}


\section{Theorems and Axioms}
\begin{axiom_stmt}[\texttt{ckks\_approximate\_correctness}]
Decryption recovers the plaintext within $2^{-\text{scaleBits}}$ precision.
\end{axiom_stmt}

\begin{axiom_stmt}[\texttt{add\_preserves\_level}]
Homomorphic addition preserves the minimum level of inputs.
\end{axiom_stmt}

\begin{axiom_stmt}[\texttt{mul\_decreases\_level}]
Homomorphic multiplication decreases level by 1.
\end{axiom_stmt}

\begin{theorem}[\texttt{depth\_needs\_levels}]
A computation of depth $d$ requires $d+1$ levels.
\end{theorem}

\begin{axiom_stmt}[\texttt{rotate\_preserves\_level}]
Slot rotation (for matrix operations) consumes no levels.
\end{axiom_stmt}

\begin{axiom_stmt}[\texttt{bootstrap\_restores\_levels}]
Bootstrapping restores levels to $L-1$.
\end{axiom_stmt}

\begin{theorem}[\texttt{transformer\_feasible}]
A transformer layer ($\sim$15 multiplications) is feasible with $\geq 15$ levels.
\end{theorem}

\begin{theorem}[\texttt{batch\_throughput}]
SIMD batching: $N/2 = 32{,}768$ values processed per ciphertext operation (for $N=2^{16}$).
\end{theorem}


\section{Lean Source}
\lstinputlisting[language=Lean]{../../formal/lean/Crypto/FHE/CKKS.lean}

\section{References}
\noindent Source code: \url{https://github.com/luxfi/proofs/blob/main/lean/Crypto/FHE/CKKS.lean}\\
\noindent White papers: \url{https://papers.lux.network}

\noindent Related white papers:
\begin{itemize}
  \item \texttt{lux-tee-computing-mesh.tex}
\end{itemize}

\end{document}
